\section{Conclusion}
\label{sec:conclusion}

The federal government is exposed, as holder, guarantor or debtor, on about
\result{SovereignUnion} percent of the claims paid directly from wages, and holds about
\result{AIFederalShareRounded} percent of the claims on AI capital; about
\result{DebtOnlyDirect} percent of all US debt, public and private, is serviced directly from
wages, government debt entering that measure through the labor-linked share of the receipts that
service it rather than because a household owes it, and the state stands behind most of it.
Displacement is a fiscal event before it is a financial one: at the one level fully inside the
observed data the federal government bears roughly four fifths of the first-round loss while the
banking system barely moves, and whether the budget absorbs that turns on which capital tax rate
reaches AI profits, since the condition needs \result{RequiredEasy} to \result{RequiredHard}
percent and the income that moves from wages to AI capital bears \result{TauKBarkai} percent once
every layer of tax is counted once. The state loses in both failure directions through different
tax bases, wage taxes if AI succeeds and capital gains and corporate taxes if it fails, so the
same capital tax rate is being asked to do two opposite jobs, and that, rather than any single
loss estimate, is what the measurement is for.

This paper is written for the people who will have to answer those questions with numbers: the
fiscal authority that has to decide whether to pre-fund, the supervisor that has to decide which
scenario to run, and the analyst who has to say how much of a balance sheet depends on wages
continuing to be paid. Everything reported here is generated from artifacts in a public
replication package, each carrying its status and its replication flag, so that a reader who
disagrees with a convention can change it and see what moves.
